\documentclass[12pt]{article}
\usepackage[T1]{fontenc}
\usepackage[utf8]{inputenc}
\usepackage{lmodern}
\usepackage{textcomp}
\usepackage{enumitem}
\usepackage{geometry}
\geometry{margin=1in}
\begin{document}
\begin{center}
Answer Key - Economics - K-12 Level Assessment 18 \
\small (Answers are ordered exactly as in the question paper.)
\end{center}

\section*{Multiple Choice}
\begin{enumerate}[label=\arabic*.]
\item \textbf{Answer:} D
\\ \textit{Options:} A) Paper producers expect lower paper prices in the months ahead.; B) The price of pencils, a complement to paper, increases.; C) Improvements in the technology used to produce paper.; D) Environmental concerns reduce the yearly amount of timber that can be harvested.
\\ \textit{Note:} Environmental concerns that limit timber harvesting reduce the availability of raw materials for paper production, leading to a decrease in supply and causing the supply curve to shift to the left.
\item \textbf{Answer:} B
\\ \textit{Options:} A) The velocity of money is equal to real GDP divided by the money supply.; B) Dollars earned today have more purchasing power than dollars earned a year from today.; C) The supply of loanable funds consists of investors.; D) The demand for loanable funds consists of savers.
\\ \textit{Note:} Dollars earned today have more purchasing power than dollars earned a year from today because of the time value of money, which indicates that money available now can be invested to earn interest or generate returns, making it more valuable than the same amount of money in the future due to inflation and opportunity cost.
\item \textbf{Answer:} A
\\ \textit{Options:} A) marginal utility is zero; B) marginal benefit equals marginal cost; C) consumer surplus is zero; D) distributive efficiency is achieved
\\ \textit{Note:} The total utility from sardines is maximized when marginal utility is zero because this indicates that consuming an additional sardine does not increase overall satisfaction, suggesting that the optimal quantity has been reached.
\item \textbf{Answer:} C
\\ \textit{Options:} A) The euro and other foreign currencies held by the Federal Reserve; B) Gold bars in secure locations like Fort Knox; C) The promise of the U.S. government to maintain its value; D) The value of the actual paper on which it is printed.
\\ \textit{Note:} The value of the U.S. dollar is insured by the trust and confidence in the U.S. government's ability to maintain economic stability and uphold its obligations, which underpins the currency's acceptance and value in the economy.
\item \textbf{Answer:} A
\\ \textit{Options:} A) Price equal to marginal cost; B) Relative welfare loss; C) Relatively high price; D) Relatively low quantity
\\ \textit{Note:} In a competitive market structure, firms are price takers and will set their prices equal to marginal cost to maximize profits, whereas a monopoly has the power to set prices above marginal cost, leading to higher prices and reduced output.
\end{enumerate}

\section*{True / False}
\begin{enumerate}[label=\arabic*.]
\setcounter{enumi}{5}
\item \textbf{Answer:} False
\\ \textit{Options:} A) True; B) False
\\ \textit{Note:} Expansionary monetary policy is designed to increase aggregate demand by lowering interest rates and increasing the money supply, which can lead to higher real prices and stimulate economic growth, not decrease demand or cause prices to fall.
\item \textbf{Answer:} True
\\ \textit{Options:} A) True; B) False
\\ \textit{Note:} The statement is true because during the decade from 1990 to 2000, many economies experienced significant increases in real GDP per capita due to factors such as technological advancements and globalization, leading to an overall growth rate of approximately 50 percent.
\item \textbf{Answer:} True
\\ \textit{Options:} A) True; B) False
\\ \textit{Note:} Economic costs encompass both explicit costs, which are direct, out-of-pocket expenses, and implicit costs, which represent the opportunity costs of foregone alternatives, thus making the statement true.
\item \textbf{Answer:} False
\\ \textit{Options:} A) True; B) False
\\ \textit{Note:} Increasing government spending can stimulate economic activity and increase aggregate demand, potentially leading to an increase in real GDP rather than a decrease.
\item \textbf{Answer:} True
\\ \textit{Options:} A) True; B) False
\\ \textit{Note:} The volunteer fire department provides essential fire protection services that benefit all members of the community, regardless of whether they directly contribute to its funding, illustrating the characteristics of a public good.
\end{enumerate}

\section*{Long Form Response}
\begin{enumerate}[label=\arabic*.]
\setcounter{enumi}{10}
\item \textbf{Answer:} A government can determine the socially optimal price for a natural monopoly by analyzing the intersection of the marginal cost curve and the demand curve. The socially optimal price occurs where the price consumers are willing to pay (demand) equals the additional cost of producing one more unit of the good (marginal cost). This intersection ensures that resources are allocated efficiently, as it maximizes total welfare for society by balancing consumer needs with production costs. If the price is set too high above this point, it could lead to a loss of consumer surplus and reduced access to the good, while a price set too low could result in losses for the producer and decreased supply. Therefore, finding this intersection helps the government set a price that reflects both the costs of production and the value to consumers.
\\ \textit{Note:} correct because it identifies that the socially optimal price for a natural monopoly is found where marginal cost equals demand, ensuring efficient resource allocation and maximization of total welfare by aligning consumer willingness to pay with production costs.
\item \textbf{Answer:} Government spending can significantly influence real gross domestic product (GDP) when an economy is operating below full employment. When the government increases spending without raising taxes, it injects more money into the economy, which can lead to higher demand for goods and services. This increased demand can stimulate production, encouraging businesses to hire more workers and invest in resources, which helps reduce unemployment. As more people earn wages and spend money, this creates a multiplier effect that can further boost overall economic growth. By strategically increasing spending, the government can help the economy recover and move closer to full employment, ultimately enhancing long-term growth prospects.
\\ \textit{Note:} correct because it illustrates how increased government spending stimulates demand, leading to higher production and employment, which ultimately contributes to economic growth when the economy is operating below full employment.
\end{enumerate}

\vfill
\noindent\textit{End of Answer Key}
\end{document}